%% LyX 2.5.1 created this file.  For more info, see https://www.lyx.org/.
%% Do not edit unless you really know what you are doing.
\documentclass[brazil]{article}
\usepackage[T1]{fontenc}
\usepackage[utf8]{inputenc}
\usepackage{geometry}
\geometry{verbose,tmargin=3cm,bmargin=3cm,lmargin=3cm,rmargin=2.5cm}
\usepackage{babel}
\begin{document}
\title{Dynamist materialism}
\maketitle
\begin{quote}
First of all, our version of materialism is dynamicist, as it identifies materiality with mutability (\cite{Bunge2013Materialismo}).
\end{quote}
Even before engaging with the philosophical literature, my contact with science had already led me to develop, in an unsystematic way, my own conception of the world. However, it was in ``Materialismo y Ciencia'' (\cite{Bunge2013Materialismo}) that I found some of these ideas more systematically developed by a modern thinker.

At first, I became interested in the critique of dialectics and sought out the book to read the corresponding chapter. Based on some readings in the field, I have always adopted a version of dialectical materialism in which the so-called “principles” of dialectics are understood in a more open-ended way, closer to research guidelines than to rigid principles. Its central components are the idea that matter is essentially dynamic and that every real object is either a system or part of a system, already using Bunge's own terminology. When Bunge proposes:
\begin{quote}
We propose the thesis that dialectical ontology has a plausible core enveloped in a mystical haze. The plausible core of dialectics consists of the hypotheses that (i) everything is involved in some process of change and (ii) at certain stages of every process, new qualities emerge.
\end{quote}
When we look at Hegelian dialectics, and even at the materialist derivation of his dialectics developed by some later thinkers, I tend to agree. However, in the way I have always sought to interpret it, its central core consists precisely of both the “materialist” and the “dynamic” elements. Thus, I see Bunge’s work as serving two main purposes: on the one hand, it systematizes and formalizes, in a more rigorous manner, a conception of the world that I had already adopted beforehand; on the other hand, it seeks to clean up the obscure vocabulary adopted by some thinkers who also identify with materialism, specifically the dialectical version.

I therefore regard his work as an effort to define more clearly the philosophical components of a materialist philosophy while rejecting obscure and excessively abstract language. At one point, Bunge remarks that dialectical materialism is “an embryo that can develop into a proper theory”; once again, I find myself in agreement. This is also why I do not identify as a Hegelian, and why I reject the more dogmatic formulations of dialectical materialism.

In general, I consider Bunge’s critique to be valid for the most part. I appreciate his concern with conceptual clarity and rigor. I have reservations only when the discussion moves into the realm of history and the social sciences, where I believe the argument may need to be formulated with somewhat greater care. When it comes to the natural sciences, however, I am in full agreement. In this regard, I would like to present Bunge’s characterization of a material object. First, we can see how it reflects the requirement of a world in which nothing is immutable:
\begin{quote}
We can therefore characterize a material object as an object that can exist in at least two states, such that it can transition from one to the other...
\end{quote}
And also that there are no isolated entities:
\begin{quote}
An object is real if and only if it influences or is influenced by another object, or is composed exclusively of real objects... Every real (material) object is a system or a component of a system. In other words, there are no isolated things.
\end{quote}
And these two elements are precisely those that I have emphasized as the pillars of my worldview since before. It is also interesting to bring up the following passage:
\begin{quote}
The preceding postulates, definitions, and theorems constitute the core of a materialist ontology with the following characteristics:

(a) \textbf{exact}: every concept is exact or can be made exact;

(b) \textbf{systematic}: every hypothesis belongs to a hypothetico-deductive system;

(c) \textbf{scientific}: every hypothesis is compatible with contemporary science;

(d) \textbf{dynamic}: every entity is mutable;

(e) \textbf{systemic}: every entity is a system or a component of some system;

(f) \textbf{emergentist}: every system possesses properties that its components do not possess;

(g) \textbf{evolutionary}: every original emergence is a stage in some evolutionary process.
\end{quote}
Overall, I find this set of characteristics appealing, but there are two points that I would like to draw particular attention to. First, regarding item (a), I do not agree that it is possible to work solely with exact concepts. The best I can conceive of doing in this regard is to make use of mathematics. However, even setting aside the question of whether it is possible to work with mathematics in an entirely exact manner, I do not believe it is possible to develop an ontology subject to such a restriction. The intrinsic ambiguity of human language involves a broad range of issues in linguistics and philosophy that, in my view, need to be taken into consideration\footnote{See, for example (\cite{Wittgenstein1953PhilosophicalInvestigations}).}.

We can interpret this characteristic in a less rigorous way, as “every concept strives to become as exact as possible”---that is, adopting as a general guiding principle the aim of developing an ontology that is always as clear and straightforward as possible. However, here we are already admitting a lower degree of exactness in the very characteristics that define our ontology.

Finally, the second point I would like to draw attention to is item (d). I would simply like to emphasize that not only is every entity mutable in the sense that it is capable of changing, but also that it necessarily cannot remain unchanged---that is, everything is in a state of constant transformation. With these observations, I believe that Bunge’s general ontology is compatible with my own, with the added benefit of being better organized and articulated.
\begin{quotation}
\bibliographystyle{plain}
\bibliography{ref}
\end{quotation}

\end{document}
